\documentclass[a4paper, 12pt]{article}
\usepackage[utf8]{inputenc}
\usepackage[english]{babel}
\usepackage{amsmath, amsfonts, amssymb, amsthm, mathtools}
\usepackage{pythonhighlight}
\usepackage{listings}
\usepackage{float}
\usepackage{makeidx}
\usepackage{siunitx}
\usepackage{array}
\usepackage[top=1in]{geometry}
\usepackage{caption}
\usepackage{subcaption}
\usepackage{graphicx}
\usepackage{siunitx}

\graphicspath{ {C:/Users/skief/Documents/UiT/semester8/space physics/progrTasks/2and3ionization/}}

\begin{document}
\author{Selma Kiefersauer}
\title{Report on Programming Task 2}
\maketitle

\section{Optical depth}

The equation for optical depth for vertical incidence:
\begin{equation}
\tau(z,\lambda, 0) = \sum_{j} \sigma_{j}^{a}(\lambda) \int_{z}^{\infty} n_j(z') \,dz'
\end{equation}
With adding the geometry for non-vertical incidence this is for some angles up to about 60\si{\degree}:
\begin{equation}
\tau(z,\lambda, \chi) =\frac{1}{cos \chi} \sum_{j} \sigma_{j}^{a}(\lambda) \int_{z}^{\infty} n_j(z') \,dz'
\end{equation}
This gets more complicated for even higher incidence angles up to 90\si{\degree}. 
\begin{equation}
\tau(z,\lambda, \chi) = \sum_{j} \sigma_{j}^{a}(\lambda) \int_{z}^{\infty} n_j(z') \left[1 -  \left( \frac{R_E + z_0}{R_E + z'} \right)^2 sin^2 \chi_0 \right]^{-1/2}  \,dz'
\end{equation}
For my calculations I used equation (2), as this reduces to equation(1) for vertical incidence anyways as cos(0)=1. 

But first the cross sections had to be interpolated to the wavelength resolution in the irradiance file. I used np.interpolate as this gave the best fiting results. But I get 0 for all wavelength after 100nm as there is no more data in the cross section file. 
\begin{figure}[H]
	\centering
	\caption{crosssections for different wavelength}
	\includegraphics[width=0.7\textwidth]{crosssec.png}
\end{figure}

Altitude variations of the thermospheric densities, based on the MSIS data file:
\begin{figure}[H]
	\centering
	\caption{}
	\includegraphics[width=0.7\textwidth]{n.png}
\end{figure}

EUV fluxes for different heights and wavelengths can be calculated as:
\begin{equation}
I(z,\lambda) = I^{\infty}(\lambda)exp(-\tau(z,\lambda,\chi))
\end{equation}
$I^{\infty}(\lambda)$ is the photon flux at infinity, or at the top of the thermosphere.

\begin{figure}[H]
	\centering
	\caption{Optical depth and EUV fluxes for different angles of incidence.}
	\begin{subfigure}[b]{0.49\textwidth}
		\includegraphics[width=\textwidth]{optdepth0}
	\end{subfigure}
	\begin{subfigure}[b]{0.49\textwidth}
		\includegraphics[width=\textwidth]{optdepth15}
	\end{subfigure}
	\begin{subfigure}[b]{0.49\textwidth}
		\includegraphics[width=\textwidth]{optdepth30}
	\end{subfigure}
	\begin{subfigure}[b]{0.49\textwidth}
		\includegraphics[width=\textwidth]{optdepth45}
	\end{subfigure}
	\begin{subfigure}[b]{0.49\textwidth}
		\includegraphics[width=\textwidth]{optdepth60}
	\end{subfigure}
		\begin{subfigure}[b]{0.49\textwidth}
		\includegraphics[width=\textwidth]{optdepth75}
	\end{subfigure}
\end{figure}

Python code for the calculations:
\begin{python}
#relevant angles for Xi
Xi = [0, 15, 30, 45, 60] #[degrees]

#optical depth for different angles of Xi, returns a list over all wavelength present in input data
#this should work for angles up to 60 degrees
def Tau(n_profiles, nr_of_species, altitude, angle):
    tau = 0
    for i in range(nr_of_species):
        integral = 0
        for z in range( altitude, len(n_data.iloc[:,0]) ):
            integral += n_profiles.iloc[z,i+1] * 1e6 * 1000
        tau += np.dot( cs[i] , integral )
    return tau * 1/np.cos( angle * (np.pi/180) )

#calculate optical depth for heights in the thermosphere (90-600km)
#and for different angles of incidence
optical_depth = [] #in [m]
for idx,X in enumerate(Xi):
    T = []
    for m in range(90,600):
        T.append( Tau(n_data, 3, m, X) )
    optical_depth.append(T)

#Irradiance !!what does it mean "/nm" in file?! need to convert units?!
def Irradiance(opt_depth):
    return np.array(I_data["irradiance"]) * np.exp(-np.array(opt_depth))


#make plots of optical depth (logarithmic scale for the optical depth?)
xcoord = I_data["wavelength"]
ycoord = n_data.iloc[91:,0]
for key,ele in enumerate(optical_depth):
    irradiance = Irradiance(ele)
    fig, axs = plt.subplots(2,1)
    plt.suptitle(f"solar zenith angle {Xi[key]} degrees")
    plot1 = axs[0].pcolormesh(xcoord, ycoord, ele, norm=mcolors.LogNorm(vmin = 1e-2, vmax= 1e3), cmap="plasma")
    plot2 = axs[1].pcolormesh(xcoord, ycoord, irradiance, norm=mcolors.LogNorm(vmin = 1e-7, vmax = 1e-2), cmap="plasma")
    cbar1 = fig.colorbar(plot1)
    cbar2 = fig.colorbar(plot2)
    cbar1.set_label("optical depth [m]")
    cbar2.set_label("Irradiance")
    axs[0].set_xlabel("wavelength [nm]")
    axs[0].set_ylabel("height [km]")
    axs[1].set_xlabel("wavelength [nm]")
    axs[1].set_ylabel("height [km]")
    plt.tight_layout()
    plt.show()
\end{python}
\end{document}